% ============================================
% Python API Reference Guide
% Practical applications using the ruopus library from Python
% ============================================

\lstdefinelanguage{Python}{
  keywords={and,as,assert,async,await,break,class,continue,def,del,elif,else,
    except,finally,for,from,global,if,import,in,is,lambda,nonlocal,not,or,
  pass,raise,return,try,while,with,yield,None,True,False,self},
  morekeywords=[2]{int,float,bool,bytes,str,list,dict,tuple,set},
  sensitive=true,
  morecomment=[l]{\#},
  morestring=[b]",
  morestring=[b]',
}

\section{Python API Reference}\label{sec:python_guide}

\subsection{Introduction}\label{subsec:python_intro}

\texttt{ruopus} ships a Python extension module built with
\href{https://pyo3.rs}{PyO3} and \href{https://numpy.org}{NumPy} interop
(\texttt{rust-numpy}). It mirrors the Rust public surface described in
\cref{sec:opus_layer} almost one-to-one: a stateful \code{OpusEncoder} and
\code{OpusDecoder}, the packet enums (\code{Mode}, \code{Bandwidth},
\code{FrameSize}, \code{Application}, \code{Signal}), an exception hierarchy
rooted at \code{OpusError}, and a \code{ruopus.lowlevel} submodule exposing
the SILK and CELT layers directly.

\begin{concept}{Zero-Copy PCM}
  PCM crosses the Rust/Python boundary as NumPy \code{float32} (or
  \code{int16}) arrays shaped \code{(frames, channels)}, with no extra copy in
  either direction: decode output is moved out of Rust's \code{Vec} directly
  into the array's buffer, and encode input is borrowed from the array via the
  buffer protocol. Every codec call releases the GIL, so decode throughput
  scales across threads exactly as the pure-Rust benchmarks in
  \cref{sec:optimizations} do.
\end{concept}

Fully type-hinted stubs (\code{.pyi}, generated from the compiled extension by
PyO3's \code{experimental-inspect} and checked with \code{pyright}) and one
set of docstrings (sourced from the Rust \code{///} comments and served to
  both IDE tooling and the Sphinx-built HTML guide at
\url{https://jmg049.github.io/ruopus/}) ship in every wheel.

\needspace{6\baselineskip}
\subsection{Installation}\label{subsec:python_install}

\begin{lstlisting}[language=bash]
pip install ruopus
\end{lstlisting}

Pre-built wheels cover CPython 3.9 through 3.13 on Linux, macOS, and Windows;
no Rust toolchain is required. NumPy 1.22+ is the only runtime dependency.
Building from source (\code{maturin build --features python}) needs a
\Q{1.85}+ Rust toolchain, matching the crate's MSRV.

\needspace{8\baselineskip}
\subsection{Quickstart: Encode and Decode}\label{subsec:python_quickstart}

The encoder and decoder are both stateful classes; construct one of each per
stream and feed them consecutive frames so cross-frame state (MDCT overlap,
mode hysteresis, concealment history) stays continuous.

\begin{lstlisting}[language=Python]
import numpy as np
import ruopus

# Stereo encoder at 64 kbps (default: fullband, complexity 10, auto mode)
enc = ruopus.OpusEncoder(2, bitrate=64_000)
dec = ruopus.OpusDecoder(2)

# 20 ms stereo frame at 48 kHz -> 960 samples per channel
frame = np.zeros((960, 2), dtype=np.float32)

packet = enc.encode_auto(frame)        # bytes: the encoded Opus packet
pcm    = dec.decode_packet(packet)     # (960, 2) float32 in [-1, 1]
\end{lstlisting}

Encoder input accepts a 1-D interleaved array or a 2-D \code{(frames,
channels)} array, both without extra copies. Mono streams use
\code{channels=1}; decoded PCM is then \code{(frames, 1)}. For callers that
need integer samples, \code{OpusDecoder}'s \code{decode\_packet\_i16} method
returns \code{int16} PCM instead, scaled and rounded exactly as
\code{opus\_demo} does.

Every encoder property (\code{bitrate}, \code{complexity}, \code{dtx},
  \code{bandwidth}, \code{signal}, \code{application}, \code{inband\_fec},
\code{packet\_loss\_perc}, \code{vbr}, \code{force\_channels}) is writable at
runtime and takes effect on the next frame -- there is no need to rebuild the
encoder to change quality or bitrate mid-stream.

\needspace{8\baselineskip}
\subsection{Choosing a Codec Mode}\label{subsec:python_modes}

\cref{sec:opus_layer} describes SILK, CELT, and hybrid mode at the bitstream
level; the Python encoder exposes one method per mode plus an automatic
dispatcher:

\begin{lstlisting}[language=Python]
enc.encode_auto(frame)     # per-frame SILK / hybrid / CELT decision
enc.encode(frame)          # force CELT-only
enc.encode_silk(frame)     # force SILK-only (speech, 10/20/40/60 ms)
enc.encode_hybrid(frame)   # force hybrid (10/20 ms, SWB/FB only)
\end{lstlisting}

\code{encode\_auto} is the right default for most applications: it runs the
encoder's signal analysis (\cref{subsec:encode_auto}) every frame and picks a
mode from the \code{signal} hint, the \code{application} profile, the target
bitrate, and cross-frame hysteresis. The \code{Signal} and \code{Application}
enums bias that decision explicitly:

\begin{lstlisting}[language=Python]
# Phone-quality voice
enc_voice = ruopus.OpusEncoder(
    1, bitrate=8_000,
    application=ruopus.Application.Voip,
    signal=ruopus.Signal.Voice,
)

# High-quality music
enc_music = ruopus.OpusEncoder(
    2, bitrate=128_000,
    application=ruopus.Application.Audio,
    signal=ruopus.Signal.Music,
)

# Low-latency game voice (CELT only, never SILK)
enc_ld = ruopus.OpusEncoder(
    1, bitrate=32_000,
    application=ruopus.Application.RestrictedLowDelay,
)
\end{lstlisting}

After encoding, \code{ruopus.Packet(packet).toc} reports which mode,
bandwidth, and frame size the encoder actually chose
(\cref{subsec:python_packets}).

\needspace{6\baselineskip}
\subsection{Frame Sizes}\label{subsec:python_framesizes}

Opus's six frame durations (\cref{tab:frame_sizes_py}) are always expressed in
samples \emph{per channel} at 48~kHz, regardless of which mode encodes them.
\code{FrameSize} carries both directions of that conversion:

\begin{lstlisting}[language=Python]
fs = ruopus.FrameSize.Ms20
print(fs.tenth_ms, fs.samples_per_channel_48k)   # 200, 960
\end{lstlisting}

\begin{table}[h]
  \centering
  \begin{tabular}{lll}
    \toprule
    Duration & Samples @ 48 kHz & Modes \\
    \midrule
    2.5 ms & 120  & CELT only \\
    5 ms   & 240  & CELT only \\
    10 ms  & 480  & SILK, hybrid, CELT \\
    20 ms  & 960  & SILK, hybrid, CELT \\
    40 ms  & 1920 & SILK only \\
    60 ms  & 2880 & SILK only \\
    \bottomrule
  \end{tabular}
  \caption{Valid frame durations and the modes that accept each.}
  \label{tab:frame_sizes_py}
\end{table}

20~ms is the only duration every mode accepts and the standard default.
Passing a duration one of \code{encode}, \code{encode\_silk}, or
\code{encode\_hybrid} does not support raises \code{EncodeError}.
\code{enc.lookahead} reports the
encoder's algorithmic delay in samples (120 for the default fullband CELT
configuration, i.e.\ the MDCT overlap of \cref{subsec:mdct}) -- skip that many
leading output samples when aligning a round-trip for testing.

\needspace{8\baselineskip}
\subsection{Ogg Opus Files}\label{subsec:python_ogg}

\cref{sec:ogg} describes the Ogg container at the bitstream level;
\code{encode\_ogg\_opus}/\code{decode\_ogg\_opus} handle the entire
round-trip -- header, pages, pre-skip, output gain -- in one call each:

\begin{lstlisting}[language=Python]
import numpy as np
import ruopus

sr = 48_000
t = np.linspace(0, 3, sr * 3, dtype=np.float32)
pcm = np.column_stack([
    np.sin(2 * np.pi * 440 * t),   # left:  A4
    np.sin(2 * np.pi * 880 * t),   # right: A5
])

ogg = ruopus.encode_ogg_opus(pcm, channels=2, bitrate=128_000)
with open("output.opus", "wb") as f:
    f.write(ogg)

with open("output.opus", "rb") as f:
    decoded_pcm, head = ruopus.decode_ogg_opus(f.read())

print(f"{head.channel_count}-ch, {decoded_pcm.shape[0] / sr:.2f}s")
\end{lstlisting}

Decoding applies the \code{OpusHead} pre-skip (\cref{subsec:granule}), trims
trailing granule padding, and applies the header's Q7.8 output gain, so the
returned PCM is ready to play or save as-is. \code{OpusHead} exposes every
identification-header field (\code{version}, \code{channel\_count},
  \code{pre\_skip}, \code{input\_sample\_rate}, \code{output\_gain\_q8}, and the
channel-mapping fields for non-family-0 streams) and round-trips back to bytes
via \code{to\_bytes()}. Only channel-mapping family 0 (mono/stereo) is
supported by \code{decode\_ogg\_opus}; surround files need the multistream
decoder (\cref{subsec:python_multistream}).

\needspace{8\baselineskip}
\subsection{Packet Loss and Forward Error Correction}\label{subsec:python_fec}

\cref{subsec:plc,subsec:fec} describe concealment and in-band FEC at the
bitstream level. The decoder exposes both directly:

\begin{lstlisting}[language=Python]
import ruopus

dec = ruopus.OpusDecoder(2)
FRAME = 960  # 20 ms at 48 kHz

for pkt in stream:
    if pkt is None:
        pcm = dec.decode_lost(frame_size=FRAME)        # concealment
    else:
        pcm = dec.decode_packet(pkt)

# Recover a lost packet from in-band FEC carried by its successor
recovered = dec.decode_fec(next_pkt, frame_size=FRAME)
\end{lstlisting}

\code{decode\_lost} extrapolates the last pitch period after a CELT/hybrid
packet; a lost frame following SILK fades to silence (SILK's own PLC is not
yet ported -- \cref{subsec:silk_plc}). \code{decode\_fec} recovers a lost
frame from the low-bitrate redundant (LBRR) copy the \emph{next} packet
carries, falling back to plain concealment when no usable FEC data is
present. Enable FEC generation on the encoder with two properties:

\begin{lstlisting}[language=Python]
enc = ruopus.OpusEncoder(
    1, bitrate=24_000,
    application=ruopus.Application.Voip,
    signal=ruopus.Signal.Voice,
    inband_fec=True,
    packet_loss_perc=10,   # hint: expect 10% loss
)
\end{lstlisting}

\code{packet\_loss\_perc} (clamped 0--100) biases the encoder toward
loss-robust coding -- more redundancy at the cost of raw efficiency -- and is
only effective on SILK/hybrid packets, so it has no effect under
\code{Application.RestrictedLowDelay}.

\needspace{8\baselineskip}
\subsection{Multistream and Surround Decoding}\label{subsec:python_multistream}

\cref{sec:multistream} describes multistream demultiplexing at the bitstream
level. \code{MultistreamDecoder} decodes \code{streams} elementary Opus
streams per packet (the first \code{coupled} as stereo, the rest mono) and
routes the resulting channels to an output layout via \code{mapping}, where
the sentinel \code{255} produces a silent channel:

\begin{lstlisting}[language=Python]
import ruopus

# 5.1 surround: 4 elementary streams, 2 stereo-coupled
dec = ruopus.MultistreamDecoder(
    streams=4,
    coupled=2,
    mapping=[0, 4, 1, 2, 3, 5],   # L R C LFE Ls Rs
)
pcm = dec.decode_packet(raw_packet)   # (frames, 6) float32
\end{lstlisting}

Each packet's elementary streams are self-delimited (\cref{subsec:self_delimit})
except the last, which is standard-framed; \code{decode\_packet} handles that
transparently. Mono/stereo Ogg files (channel-mapping family 0) never use
multistream framing, so \code{decode\_ogg\_opus} handles them without this
class.

\needspace{8\baselineskip}
\subsection{Packet Introspection}\label{subsec:python_packets}

\code{Toc} and \code{Packet} parse Opus framing (\cref{sec:packet}) without
decoding audio -- useful for logging, monitoring, or building test vectors:

\begin{lstlisting}[language=Python]
import ruopus

pkt = ruopus.Packet(raw_bytes)
print(f"mode={pkt.toc.mode}, bandwidth={pkt.toc.bandwidth}, frames={len(pkt)}")

for frame in pkt.frames:      # each a `bytes` object; may be empty (DTX)
    ...
first_frame = pkt[0]           # Packet is a read-only sequence of frames
\end{lstlisting}

\code{Toc} exposes the raw byte plus its three decoded fields (\code{mode},
  \code{bandwidth}, \code{frame\_size}, \code{channels}, \code{stereo},
\code{config}, \code{frame\_count\_code}) and is hashable, so TOC
distributions across a stream can be tallied with a \code{Counter}.
\code{Packet.parse} validates the R1--R7 framing rules
(\cref{subsec:validity}) and raises \code{PacketError} on a malformed packet;
\code{Packet.parse\_self\_delimited} parses the self-delimiting form used
inside multistream payloads and returns the bytes consumed alongside the
packet. After a round-trip, comparing \code{enc.final\_range} against
\code{dec.final\_range} (\cref{subsec:oracle}) is the same bit-exactness
oracle used by the Rust conformance suite.

\needspace{8\baselineskip}
\subsection{The Low-Level API}\label{subsec:python_lowlevel}

\code{ruopus.lowlevel} exposes the SILK and CELT layers of
\cref{sec:celt,sec:silk} directly, below the Opus packet format, plus the LPC
analysis/synthesis pipeline underlying SILK. These are research and
custom-codec building blocks; ordinary use should prefer
\code{OpusEncoder}/\code{OpusDecoder}.

\begin{lstlisting}[language=Python]
import numpy as np
import ruopus.lowlevel as ll

# The CELT layer alone: raw frame bodies, no TOC, no Opus framing
enc = ll.CeltEncoder(channels=1, complexity=10, bitrate=64_000)
frame = np.zeros(960, dtype=np.float32)
body = enc.encode(frame)              # raw CELT frame body

# SILK operates at its own internal rate (8/12/16 kHz), not 48 kHz,
# and on int16 PCM
silk_enc = ll.SilkEncoder(fs_khz=16, nb_subfr=4, bitrate=25_000)
payload = silk_enc.encode(np.zeros(320, dtype=np.int16))  # 20 ms @ 16 kHz

# The LPC analysis pipeline that drives SILK's short-term prediction
signal = np.random.randn(320).astype(np.float32)
coeffs = ll.lpc_analysis(signal, order=16)                # Levinson-Durbin
residual = ll.lpc_residual_stateful(signal, coeffs)
lag = ll.estimate_pitch(signal, fs_khz=16)                 # long-term prediction
\end{lstlisting}

\code{SilkStereoEncoder} extends \code{SilkEncoder} with mid-side coding
(\cref{subsec:silk_stereo}); \code{DecControl} carries the SILK decoder's
per-frame control parameters. \code{CeltDecoder} does not expose a full
\code{decode} method, since CELT decoding needs the encoder's live
range-coder state (\cref{subsec:oracle}), which crosses no API boundary --
real CELT packets decode through \code{OpusDecoder} instead; the low-level
decoder covers concealment (\code{decode\_lost}) independently.

\needspace{6\baselineskip}
\subsection{Complete API Reference}\label{subsec:python_api_ref}

\begin{table}[h]
  \centering
  \begin{tabular}{ll}
    \toprule
    Name & Purpose \\
    \midrule
    \code{OpusEncoder} / \code{OpusDecoder} & Top-level Opus packet
    codec (\cref{subsec:python_quickstart}) \\
    \code{MultistreamDecoder} & Surround/multistream decode
    (\cref{subsec:python_multistream}) \\
    \code{Packet} / \code{Toc} & Packet and TOC introspection
    (\cref{subsec:python_packets}) \\
    \code{OpusHead} & Ogg Opus identification header
    (\cref{subsec:python_ogg}) \\
    \code{encode\_ogg\_opus} / \code{decode\_ogg\_opus} & Whole-file
    Ogg Opus helpers (\cref{subsec:python_ogg}) \\
    \code{Mode}, \code{Bandwidth}, \code{FrameSize} & Packet-level
    enums (\cref{sec:packet}) \\
    \code{Signal}, \code{Application} & Encoder mode/latency hints
    (\cref{subsec:python_modes}) \\
    \code{OpusError} & Base exception for every codec failure \\
    \code{PacketError}, \code{EncodeError}, \code{OggError} &
    Specific exception subclasses \\
    \code{ruopus.lowlevel} & SILK/CELT/LPC primitives
    (\cref{subsec:python_lowlevel}) \\
    \code{ruopus.version()} & The installed \texttt{ruopus} version string \\
    \bottomrule
  \end{tabular}
  \caption{The top-level \texttt{ruopus} Python surface. Every class and
    function carries a full docstring and type stub; run \code{help(ruopus.X)}
    or consult \url{https://jmg049.github.io/ruopus/} for exhaustive
  parameter-level documentation.}
  \label{tab:python_api_ref}
\end{table}

\begin{importantnote}
  This reference is deliberately a map, not the territory: exhaustive
  parameter documentation lives in the generated \code{.pyi} stubs (so IDEs
  show it inline) and the Sphinx guide. Anywhere this section and the
  installed package disagree, the package is correct -- file an issue against
  the manual.
\end{importantnote}
